\capitulo{7}{Conclusiones y líneas de trabajo futuras}

\section{Conclusiones}

El desarrollo de este Trabajo de Fin de Grado ha terminado de manera satisfactoria, logrando el cumplimiento de los objetivos propuestos desde su planificación. Con AgroSkopos, se ha materializado una plataforma web centralizada que da respuesta a las necesidades de monitorización, telemetría y control remoto en el ámbito de la agricultura de precisión. Mediante una arquitectura de baja latencia, se han satisfecho con éxito los requisitos funcionales exigidos: desde la implementación de un mapa interactivo y la gestión de misiones autónomas, hasta el despliegue de un panel de telemetría y control manual en tiempo real. Además, la adopción de metodologías ágiles y herramientas de análisis estático ha garantizado un producto final robusto, contenerizado y con altos estándares de calidad.

Desde la perspectiva de la ingeniería del \textit{software}, el proyecto ha representado un valioso ejercicio de diseño arquitectónico y resolución de problemas complejos. El principal reto superado ha sido la construcción de una infraestructura escalable y de baja latencia, orquestando con éxito la \textbf{comunicación bidireccional en tiempo real} (WebSockets), la gestión eficiente del estado en el cliente y la automatización de pruebas mediante flujos de integración continua (CI/CD). Esta robustez estructural fue precisamente la que permitió una rápida adaptación técnica ante el retraso logístico del \textit{hardware}.

En el plano formativo, este proyecto ha supuesto un cierre perfecto para mi formación en el grado. Asumir en solitario la responsabilidad integral del ciclo de vida del \textit{software}, desde la toma de requisitos hasta las pruebas automatizadas, ha permitido consolidar y poner en práctica de manera transversal los conocimientos teóricos adquiridos a lo largo de la carrera. Simultáneamente, el proyecto ha motivado un enriquecedor proceso de autoaprendizaje, fundamental para dominar y aplicar tecnologías modernas ampliamente demandadas en el sector.

En definitiva, AgroSkopos se consolida como un ecosistema tecnológico integral orientado a la digitalización del sector agrario. Su evolución desde un prototipo conceptual hasta una plataforma \textit{cloud-ready}, asíncrona y sometida a rigurosas pruebas de carga, demuestra la viabilidad de aplicar arquitecturas de \textit{software} modernas (propias de sistemas de alta transaccionalidad) al control en tiempo real de maquinaria pesada.

\section{Líneas de trabajo futuras}

El desarrollo y la posterior evaluación de AgroSkopos han puesto de manifiesto la viabilidad de la plataforma como núcleo tecnológico operativo, abriendo a su vez nuevas e interesantes vías de investigación e ingeniería para enriquecer el sistema en futuras iteraciones.

\subsection{Transición a \textit{Hardware} Real y Comunicaciones Avanzadas}

El sistema actual se apoya de forma robusta en un simulador nativo diseñado para replicar la dinámica del entorno real. El siguiente paso natural en la evolución del proyecto es sustituir este motor de \textit{software} por una \textbf{pasarela IoT} (\textit{Internet of Things}) que establezca comunicación directa con el chasis de un robot físico real. Para lograrlo, se propone el uso del estándar \textit{Robot Operating System} (ROS2) integrado con MQTT, lo que garantizaría un flujo telemétrico eficiente. 

Además, actualmente el visor de cámara emplea una simulación funcional integrada en la interfaz; la mejora propuesta consiste en implementar un servidor de \textit{streaming} nativo basado en la tecnología \textbf{WebRTC}. Esto permitirá transmitir el flujo de vídeo en vivo directamente desde las cámaras del robot hacia el panel de control del operador con latencias sub-segundo, un factor crítico e indispensable para asegurar la maniobrabilidad durante la operación manual a distancia.

Por otra parte, se plantea adaptar la arquitectura del sistema para soportar \textbf{operaciones sin conexión} en entornos agrícolas caracterizados por conectividad intermitente. La solución pasaría por dotar al cliente de la capacidad de almacenar localmente tanto la telemetría generada como el registro temporal de misiones y captura de datos. Al recuperar la señal de red, el sistema orquestará una sincronización diferida automática contra la base de datos central en PostgreSQL \cite{postgresql}, asegurando la integridad transaccional de los datos agrícolas sin riesgo de pérdida.

\subsection{Visión Artificial, Realidad Aumentada y Registro Visual}

Aprovechando el flujo continuo de vídeo en tiempo real implementado en el control remoto, el ecosistema puede beneficiarse enormemente de la integración de modelos de \textbf{\textit{Machine Learning}} (tales como la arquitectura YOLOv8). Al procesar las imágenes en el propio extremo de la red (\textbf{\textit{Edge Computing}}), la inteligencia artificial sería capaz de identificar amenazas agronómicas, como la presencia de malas hierbas o indicadores visuales de enfermedades en las plantaciones.

Mediante el uso de librerías nativas del navegador web como WebGL o Canvas API, la información detectada por la IA se proyectaría directamente sobre el vídeo simulando una interfaz de tipo HUD (\textit{Head-Up Display}). Este enfoque enriquecería sustancialmente la experiencia del operador, permitiéndole no solo observar el terreno a través de la cámara, sino también visualizar marcadores geográficos precisos, trazados dinámicos de rutas, delimitaciones de geocercas y alertas tempranas, multiplicando la conciencia situacional en tiempo real.

Complementariamente, el sistema integraría un módulo de captura y almacenamiento geolocalizado de imágenes. Esto facilitaría el trabajo de ingeniería agrícola al permitir la configuración autónoma del robot para registrar el estado visual del cultivo fotograma a fotograma en los distintos puntos de recolección de datos (\textit{waypoints}), generando en la base de datos un inventario y un rastro visual persistente asociado a las coordenadas GPS exactas.

\subsection{Planificación Óptima y Control Cinemático Avanzado}

Aunque los algoritmos de navegación actuales garantizan una cobertura de área funcional (mediante patrones estáticos en \textit{zigzag}), existe un gran margen de mejora hacia la planificación adaptativa inteligente. El sistema evolucionaría con la introducción de \textbf{algoritmos de búsqueda heurística de grafos}, tales como A* o D* Lite, los cuales se aplicarían no sobre un plano liso, sino sobre un sofisticado "mapa de costes" topográfico. La IA evaluaría en tiempo real diversas condiciones críticas del entorno (desde la orografía y el ángulo de inclinación del terreno, hasta la densidad del suelo) para trazar la ruta óptima de menor esfuerzo cinemático, reduciendo drásticamente tanto el tiempo empleado por misión como el consumo energético total.

En escenarios reales de operación, la seguridad del robot frente a imprevistos es ineludible. Por este motivo, la plataforma debe integrar en sus algoritmos el procesamiento en bruto procedente de escáneres LIDAR y sensores de proximidad ultrasonidos. La aplicación de técnicas de \textbf{evasión dinámica de obstáculos} (como los Histogramas de Campo Vectorial) facultaría al robot agrícola para sortear automáticamente maquinaria u operarios que interfirieran accidentalmente en su trayectoria predefinida, garantizando una maniobra segura y la posterior corrección automatizada hacia el objetivo original.

\subsection{Escalabilidad hacia la Gestión de Flotas (\textit{Swarm Robotics})}

AgroSkopos, en su versión fundacional, ha optimizado su capa de servicios y comunicación de eventos asumiendo un escenario operativo de un único vehículo agrícola. No obstante, la robustez de su motor de bases de datos PostgreSQL \cite{postgresql} y su canal de WebSockets soportan teóricamente arquitecturas de mayor densidad y demanda de flujo. La proyección hacia el futuro de esta herramienta radicaría en rediseñar la plataforma de control como un \textbf{administrador de flotas centralizado}, dotando a la interfaz de las capacidades de monitorización múltiple y en paralelo. Esta evolución orquestaría escenarios de cooperación avanzada, posibilitando a un único operador humano la supervisión general y la delegación de tareas concurrentes a múltiples robots agrícolas (\textit{enjambres}) operando en sincronía a lo largo de extensiones de cultivo a gran escala.

\subsection{Migración a Tipado Estricto (TypeScript)}

Por último, para dotar al producto de un estándar de mantenibilidad superior y facilitar la incorporación futura de nuevos desarrolladores al equipo, el proceso lógico siguiente consistiría en la migración progresiva del código fuente (actualmente escrito en JavaScript) hacia \textbf{TypeScript}. Adoptar un sistema de tipado estricto en el compilador permitiría detectar errores semánticos y de interfaz en etapas tempranas del desarrollo, ofreciendo un contrato robusto de tipos de datos, especialmente crítico cuando se manejan estructuras telemétricas tan complejas y densas provenientes del robot.
